\documentclass[11pt]{article}

    % --- NUESTRO FORMATO REQUERIDO ---
    \usepackage{setspace}
    \onehalfspacing

    \usepackage{fontspec}
    \setmainfont{Carlito-Regular.ttf}[
        BoldFont = Carlito-Bold.ttf,
        ItalicFont = Carlito-Italic.ttf,
        BoldItalicFont = Carlito-BoldItalic.ttf
    ]
    \setsansfont{Carlito-Regular.ttf}[
        BoldFont = Carlito-Bold.ttf,
        ItalicFont = Carlito-Italic.ttf,
        BoldItalicFont = Carlito-BoldItalic.ttf
    ]
    \renewcommand{\familydefault}{\sfdefault}

    \usepackage{hyperref}
    \hypersetup{
        pdftitle={Actividad 2},
        pdfauthor={Héctor Lobato Silva}
    }
    % ---------------------------------

    \usepackage[breakable]{tcolorbox}
    \usepackage{parskip} % Stop auto-indenting (to mimic markdown behaviour)
    

    % Basic figure setup, for now with no caption control since it's done
    % automatically by Pandoc (which extracts ![](path) syntax from Markdown).
    \usepackage{graphicx}
    % Keep aspect ratio if custom image width or height is specified
    \setkeys{Gin}{keepaspectratio}
    % Maintain compatibility with old templates. Remove in nbconvert 6.0
    \let\Oldincludegraphics\includegraphics
    % Ensure that by default, figures have no caption (until we provide a
    % proper Figure object with a Caption API and a way to capture that
    % in the conversion process - todo).
    \usepackage{caption}
    \DeclareCaptionFormat{nocaption}{}
    \captionsetup{format=nocaption,aboveskip=0pt,belowskip=0pt}

    \usepackage{float}
    \floatplacement{figure}{H} % forces figures to be placed at the correct location
    \usepackage{xcolor} % Allow colors to be defined
    \usepackage{enumerate} % Needed for markdown enumerations to work
    \usepackage{geometry} % Used to adjust the document margins
    \usepackage{amsmath} % Equations
    \usepackage{amssymb} % Equations
    \usepackage{textcomp} % defines textquotesingle
    % Hack from http://tex.stackexchange.com/a/47451/13684:
    \AtBeginDocument{%
        \def\PYZsq{\textquotesingle}% Upright quotes in Pygmentized code
    }
    \usepackage{upquote} % Upright quotes for verbatim code
    \usepackage{eurosym} % defines \euro

    \usepackage{iftex}
    \ifPDFTeX
        \usepackage[T1]{fontenc}
        \IfFileExists{alphabeta.sty}{
              \usepackage{alphabeta}
          }{
              \usepackage[mathletters]{ucs}
              \usepackage[utf8x]{inputenc}
          }
    \else
        \usepackage{fontspec}
        \usepackage{unicode-math}
    \fi

    \usepackage{fancyvrb} % verbatim replacement that allows latex
    \usepackage{grffile} % extends the file name processing of package graphics
                         % to support a larger range
    \makeatletter % fix for old versions of grffile with XeLaTeX
    \@ifpackagelater{grffile}{2019/11/01}
    {
      % Do nothing on new versions
    }
    {
      \def\Gread@@xetex#1{%
        \IfFileExists{"\Gin@base".bb}%
        {\Gread@eps{\Gin@base.bb}}%
        {\Gread@@xetex@aux#1}%
      }
    }
    \makeatother
    \usepackage[Export]{adjustbox} % Used to constrain images to a maximum size
    \adjustboxset{max size={0.9\linewidth}{0.9\paperheight}}

    % The hyperref package gives us a pdf with properly built
    % internal navigation ('pdf bookmarks' for the table of contents,
    % internal cross-reference links, web links for URLs, etc.)
    \usepackage{hyperref}
    % The default LaTeX title has an obnoxious amount of whitespace. By default,
    % titling removes some of it. It also provides customization options.
    \usepackage{titling}
    \usepackage{longtable} % longtable support required by pandoc >1.10
    \usepackage{booktabs}  % table support for pandoc > 1.12.2
    \usepackage{array}     % table support for pandoc >= 2.11.3
    \usepackage{calc}      % table minipage width calculation for pandoc >= 2.11.1
    \usepackage[inline]{enumitem} % IRkernel/repr support (it uses the enumerate* environment)
    \usepackage[normalem]{ulem} % ulem is needed to support strikethroughs (\sout)
                                % normalem makes italics be italics, not underlines
    \usepackage{soul}      % strikethrough (\st) support for pandoc >= 3.0.0
    \usepackage{mathrsfs}
    

    
    % Colors for the hyperref package
    \definecolor{urlcolor}{rgb}{0,.145,.698}
    \definecolor{linkcolor}{rgb}{.71,0.21,0.01}
    \definecolor{citecolor}{rgb}{.12,.54,.11}

    % ANSI colors
    \definecolor{ansi-black}{HTML}{3E424D}
    \definecolor{ansi-black-intense}{HTML}{282C36}
    \definecolor{ansi-red}{HTML}{E75C58}
    \definecolor{ansi-red-intense}{HTML}{B22B31}
    \definecolor{ansi-green}{HTML}{00A250}
    \definecolor{ansi-green-intense}{HTML}{007427}
    \definecolor{ansi-yellow}{HTML}{DDB62B}
    \definecolor{ansi-yellow-intense}{HTML}{B27D12}
    \definecolor{ansi-blue}{HTML}{208FFB}
    \definecolor{ansi-blue-intense}{HTML}{0065CA}
    \definecolor{ansi-magenta}{HTML}{D160C4}
    \definecolor{ansi-magenta-intense}{HTML}{A03196}
    \definecolor{ansi-cyan}{HTML}{60C6C8}
    \definecolor{ansi-cyan-intense}{HTML}{258F8F}
    \definecolor{ansi-white}{HTML}{C5C1B4}
    \definecolor{ansi-white-intense}{HTML}{A1A6B2}
    \definecolor{ansi-default-inverse-fg}{HTML}{FFFFFF}
    \definecolor{ansi-default-inverse-bg}{HTML}{000000}

    % common color for the border for error outputs.
    \definecolor{outerrorbackground}{HTML}{FFDFDF}

    % commands and environments needed by pandoc snippets
    % extracted from the output of `pandoc -s`
    \providecommand{\tightlist}{%
      \setlength{\itemsep}{0pt}\setlength{\parskip}{0pt}}
    \DefineVerbatimEnvironment{Highlighting}{Verbatim}{commandchars=\\\{\}}
    % Add ',fontsize=\small' for more characters per line
    \newenvironment{Shaded}{}{}
    \newcommand{\KeywordTok}[1]{\textcolor[rgb]{0.00,0.44,0.13}{\textbf{{#1}}}}
    \newcommand{\DataTypeTok}[1]{\textcolor[rgb]{0.56,0.13,0.00}{{#1}}}
    \newcommand{\DecValTok}[1]{\textcolor[rgb]{0.25,0.63,0.44}{{#1}}}
    \newcommand{\BaseNTok}[1]{\textcolor[rgb]{0.25,0.63,0.44}{{#1}}}
    \newcommand{\FloatTok}[1]{\textcolor[rgb]{0.25,0.63,0.44}{{#1}}}
    \newcommand{\CharTok}[1]{\textcolor[rgb]{0.25,0.44,0.63}{{#1}}}
    \newcommand{\StringTok}[1]{\textcolor[rgb]{0.25,0.44,0.63}{{#1}}}
    \newcommand{\CommentTok}[1]{\textcolor[rgb]{0.38,0.63,0.69}{\textit{{#1}}}}
    \newcommand{\OtherTok}[1]{\textcolor[rgb]{0.00,0.44,0.13}{{#1}}}
    \newcommand{\AlertTok}[1]{\textcolor[rgb]{1.00,0.00,0.00}{\textbf{{#1}}}}
    \newcommand{\FunctionTok}[1]{\textcolor[rgb]{0.02,0.16,0.49}{{#1}}}
    \newcommand{\RegionMarkerTok}[1]{{#1}}
    \newcommand{\ErrorTok}[1]{\textcolor[rgb]{1.00,0.00,0.00}{\textbf{{#1}}}}
    \newcommand{\NormalTok}[1]{{#1}}

    % Additional commands for more recent versions of Pandoc
    \newcommand{\ConstantTok}[1]{\textcolor[rgb]{0.53,0.00,0.00}{{#1}}}
    \newcommand{\SpecialCharTok}[1]{\textcolor[rgb]{0.25,0.44,0.63}{{#1}}}
    \newcommand{\VerbatimStringTok}[1]{\textcolor[rgb]{0.25,0.44,0.63}{{#1}}}
    \newcommand{\SpecialStringTok}[1]{\textcolor[rgb]{0.73,0.40,0.53}{{#1}}}
    \newcommand{\ImportTok}[1]{{#1}}
    \newcommand{\DocumentationTok}[1]{\textcolor[rgb]{0.73,0.13,0.13}{\textit{{#1}}}}
    \newcommand{\AnnotationTok}[1]{\textcolor[rgb]{0.38,0.63,0.69}{\textbf{\textit{{#1}}}}}
    \newcommand{\CommentVarTok}[1]{\textcolor[rgb]{0.38,0.63,0.69}{\textbf{\textit{{#1}}}}}
    \newcommand{\VariableTok}[1]{\textcolor[rgb]{0.10,0.09,0.49}{{#1}}}
    \newcommand{\ControlFlowTok}[1]{\textcolor[rgb]{0.00,0.44,0.13}{\textbf{{#1}}}}
    \newcommand{\OperatorTok}[1]{\textcolor[rgb]{0.40,0.40,0.40}{{#1}}}
    \newcommand{\BuiltInTok}[1]{{#1}}
    \newcommand{\ExtensionTok}[1]{{#1}}
    \newcommand{\PreprocessorTok}[1]{\textcolor[rgb]{0.74,0.48,0.00}{{#1}}}
    \newcommand{\AttributeTok}[1]{\textcolor[rgb]{0.49,0.56,0.16}{{#1}}}
    \newcommand{\InformationTok}[1]{\textcolor[rgb]{0.38,0.63,0.69}{\textbf{\textit{{#1}}}}}
    \newcommand{\WarningTok}[1]{\textcolor[rgb]{0.38,0.63,0.69}{\textbf{\textit{{#1}}}}}
    \makeatletter
    \newsavebox\pandoc@box
    \newcommand*\pandocbounded[1]{%
      \sbox\pandoc@box{#1}%
      % scaling factors for width and height
      \Gscale@div\@tempa\textheight{\dimexpr\ht\pandoc@box+\dp\pandoc@box\relax}%
      \Gscale@div\@tempb\linewidth{\wd\pandoc@box}%
      % select the smaller of both
      \ifdim\@tempb\p@<\@tempa\p@
        \let\@tempa\@tempb
      \fi
      % scaling accordingly (\@tempa < 1)
      \ifdim\@tempa\p@<\p@
        \scalebox{\@tempa}{\usebox\pandoc@box}%
      % scaling not needed, use as it is
      \else
        \usebox{\pandoc@box}%
      \fi
    }
    \makeatother

    % Define a nice break command that doesn't care if a line doesn't already
    % exist.
    \def\br{\hspace*{\fill} \\* }
    % Math Jax compatibility definitions
    \def\gt{>}
    \def\lt{<}
    \let\Oldtex\TeX
    \let\Oldlatex\LaTeX
    \renewcommand{\TeX}{\textrm{\Oldtex}}
    \renewcommand{\LaTeX}{\textrm{\Oldlatex}}
    % Document parameters
    % Document title
    \title{Actividad 2}
    
    
    
    
    
    
    
% Pygments definitions
\makeatletter
\def\PY@reset{\let\PY@it=\relax \let\PY@bf=\relax%
    \let\PY@ul=\relax \let\PY@tc=\relax%
    \let\PY@bc=\relax \let\PY@ff=\relax}
\def\PY@tok#1{\csname PY@tok@#1\endcsname}
\def\PY@toks#1+{\ifx\relax#1\empty\else%
    \PY@tok{#1}\expandafter\PY@toks\fi}
\def\PY@do#1{\PY@bc{\PY@tc{\PY@ul{%
    \PY@it{\PY@bf{\PY@ff{#1}}}}}}}
\def\PY#1#2{\PY@reset\PY@toks#1+\relax+\PY@do{#2}}

\@namedef{PY@tok@w}{\def\PY@tc##1{\textcolor[rgb]{0.73,0.73,0.73}{##1}}}
\@namedef{PY@tok@c}{\let\PY@it=\textit\def\PY@tc##1{\textcolor[rgb]{0.24,0.48,0.48}{##1}}}
\@namedef{PY@tok@cp}{\def\PY@tc##1{\textcolor[rgb]{0.61,0.40,0.00}{##1}}}
\@namedef{PY@tok@k}{\let\PY@bf=\textbf\def\PY@tc##1{\textcolor[rgb]{0.00,0.50,0.00}{##1}}}
\@namedef{PY@tok@kp}{\def\PY@tc##1{\textcolor[rgb]{0.00,0.50,0.00}{##1}}}
\@namedef{PY@tok@kt}{\def\PY@tc##1{\textcolor[rgb]{0.69,0.00,0.25}{##1}}}
\@namedef{PY@tok@o}{\def\PY@tc##1{\textcolor[rgb]{0.40,0.40,0.40}{##1}}}
\@namedef{PY@tok@ow}{\let\PY@bf=\textbf\def\PY@tc##1{\textcolor[rgb]{0.67,0.13,1.00}{##1}}}
\@namedef{PY@tok@nb}{\def\PY@tc##1{\textcolor[rgb]{0.00,0.50,0.00}{##1}}}
\@namedef{PY@tok@nf}{\def\PY@tc##1{\textcolor[rgb]{0.00,0.00,1.00}{##1}}}
\@namedef{PY@tok@nc}{\let\PY@bf=\textbf\def\PY@tc##1{\textcolor[rgb]{0.00,0.00,1.00}{##1}}}
\@namedef{PY@tok@nn}{\let\PY@bf=\textbf\def\PY@tc##1{\textcolor[rgb]{0.00,0.00,1.00}{##1}}}
\@namedef{PY@tok@ne}{\let\PY@bf=\textbf\def\PY@tc##1{\textcolor[rgb]{0.80,0.25,0.22}{##1}}}
\@namedef{PY@tok@nv}{\def\PY@tc##1{\textcolor[rgb]{0.10,0.09,0.49}{##1}}}
\@namedef{PY@tok@no}{\def\PY@tc##1{\textcolor[rgb]{0.53,0.00,0.00}{##1}}}
\@namedef{PY@tok@nl}{\def\PY@tc##1{\textcolor[rgb]{0.46,0.46,0.00}{##1}}}
\@namedef{PY@tok@ni}{\let\PY@bf=\textbf\def\PY@tc##1{\textcolor[rgb]{0.44,0.44,0.44}{##1}}}
\@namedef{PY@tok@na}{\def\PY@tc##1{\textcolor[rgb]{0.41,0.47,0.13}{##1}}}
\@namedef{PY@tok@nt}{\let\PY@bf=\textbf\def\PY@tc##1{\textcolor[rgb]{0.00,0.50,0.00}{##1}}}
\@namedef{PY@tok@nd}{\def\PY@tc##1{\textcolor[rgb]{0.67,0.13,1.00}{##1}}}
\@namedef{PY@tok@s}{\def\PY@tc##1{\textcolor[rgb]{0.73,0.13,0.13}{##1}}}
\@namedef{PY@tok@sd}{\let\PY@it=\textit\def\PY@tc##1{\textcolor[rgb]{0.73,0.13,0.13}{##1}}}
\@namedef{PY@tok@si}{\let\PY@bf=\textbf\def\PY@tc##1{\textcolor[rgb]{0.64,0.35,0.47}{##1}}}
\@namedef{PY@tok@se}{\let\PY@bf=\textbf\def\PY@tc##1{\textcolor[rgb]{0.67,0.36,0.12}{##1}}}
\@namedef{PY@tok@sr}{\def\PY@tc##1{\textcolor[rgb]{0.64,0.35,0.47}{##1}}}
\@namedef{PY@tok@ss}{\def\PY@tc##1{\textcolor[rgb]{0.10,0.09,0.49}{##1}}}
\@namedef{PY@tok@sx}{\def\PY@tc##1{\textcolor[rgb]{0.00,0.50,0.00}{##1}}}
\@namedef{PY@tok@m}{\def\PY@tc##1{\textcolor[rgb]{0.40,0.40,0.40}{##1}}}
\@namedef{PY@tok@gh}{\let\PY@bf=\textbf\def\PY@tc##1{\textcolor[rgb]{0.00,0.00,0.50}{##1}}}
\@namedef{PY@tok@gu}{\let\PY@bf=\textbf\def\PY@tc##1{\textcolor[rgb]{0.50,0.00,0.50}{##1}}}
\@namedef{PY@tok@gd}{\def\PY@tc##1{\textcolor[rgb]{0.63,0.00,0.00}{##1}}}
\@namedef{PY@tok@gi}{\def\PY@tc##1{\textcolor[rgb]{0.00,0.52,0.00}{##1}}}
\@namedef{PY@tok@gr}{\def\PY@tc##1{\textcolor[rgb]{0.89,0.00,0.00}{##1}}}
\@namedef{PY@tok@ge}{\let\PY@it=\textit}
\@namedef{PY@tok@gs}{\let\PY@bf=\textbf}
\@namedef{PY@tok@ges}{\let\PY@bf=\textbf\let\PY@it=\textit}
\@namedef{PY@tok@gp}{\let\PY@bf=\textbf\def\PY@tc##1{\textcolor[rgb]{0.00,0.00,0.50}{##1}}}
\@namedef{PY@tok@go}{\def\PY@tc##1{\textcolor[rgb]{0.44,0.44,0.44}{##1}}}
\@namedef{PY@tok@gt}{\def\PY@tc##1{\textcolor[rgb]{0.00,0.27,0.87}{##1}}}
\@namedef{PY@tok@err}{\def\PY@bc##1{{\setlength{\fboxsep}{\string -\fboxrule}\fcolorbox[rgb]{1.00,0.00,0.00}{1,1,1}{\strut ##1}}}}
\@namedef{PY@tok@kc}{\let\PY@bf=\textbf\def\PY@tc##1{\textcolor[rgb]{0.00,0.50,0.00}{##1}}}
\@namedef{PY@tok@kd}{\let\PY@bf=\textbf\def\PY@tc##1{\textcolor[rgb]{0.00,0.50,0.00}{##1}}}
\@namedef{PY@tok@kn}{\let\PY@bf=\textbf\def\PY@tc##1{\textcolor[rgb]{0.00,0.50,0.00}{##1}}}
\@namedef{PY@tok@kr}{\let\PY@bf=\textbf\def\PY@tc##1{\textcolor[rgb]{0.00,0.50,0.00}{##1}}}
\@namedef{PY@tok@bp}{\def\PY@tc##1{\textcolor[rgb]{0.00,0.50,0.00}{##1}}}
\@namedef{PY@tok@fm}{\def\PY@tc##1{\textcolor[rgb]{0.00,0.00,1.00}{##1}}}
\@namedef{PY@tok@vc}{\def\PY@tc##1{\textcolor[rgb]{0.10,0.09,0.49}{##1}}}
\@namedef{PY@tok@vg}{\def\PY@tc##1{\textcolor[rgb]{0.10,0.09,0.49}{##1}}}
\@namedef{PY@tok@vi}{\def\PY@tc##1{\textcolor[rgb]{0.10,0.09,0.49}{##1}}}
\@namedef{PY@tok@vm}{\def\PY@tc##1{\textcolor[rgb]{0.10,0.09,0.49}{##1}}}
\@namedef{PY@tok@sa}{\def\PY@tc##1{\textcolor[rgb]{0.73,0.13,0.13}{##1}}}
\@namedef{PY@tok@sb}{\def\PY@tc##1{\textcolor[rgb]{0.73,0.13,0.13}{##1}}}
\@namedef{PY@tok@sc}{\def\PY@tc##1{\textcolor[rgb]{0.73,0.13,0.13}{##1}}}
\@namedef{PY@tok@dl}{\def\PY@tc##1{\textcolor[rgb]{0.73,0.13,0.13}{##1}}}
\@namedef{PY@tok@s2}{\def\PY@tc##1{\textcolor[rgb]{0.73,0.13,0.13}{##1}}}
\@namedef{PY@tok@sh}{\def\PY@tc##1{\textcolor[rgb]{0.73,0.13,0.13}{##1}}}
\@namedef{PY@tok@s1}{\def\PY@tc##1{\textcolor[rgb]{0.73,0.13,0.13}{##1}}}
\@namedef{PY@tok@mb}{\def\PY@tc##1{\textcolor[rgb]{0.40,0.40,0.40}{##1}}}
\@namedef{PY@tok@mf}{\def\PY@tc##1{\textcolor[rgb]{0.40,0.40,0.40}{##1}}}
\@namedef{PY@tok@mh}{\def\PY@tc##1{\textcolor[rgb]{0.40,0.40,0.40}{##1}}}
\@namedef{PY@tok@mi}{\def\PY@tc##1{\textcolor[rgb]{0.40,0.40,0.40}{##1}}}
\@namedef{PY@tok@il}{\def\PY@tc##1{\textcolor[rgb]{0.40,0.40,0.40}{##1}}}
\@namedef{PY@tok@mo}{\def\PY@tc##1{\textcolor[rgb]{0.40,0.40,0.40}{##1}}}
\@namedef{PY@tok@ch}{\let\PY@it=\textit\def\PY@tc##1{\textcolor[rgb]{0.24,0.48,0.48}{##1}}}
\@namedef{PY@tok@cm}{\let\PY@it=\textit\def\PY@tc##1{\textcolor[rgb]{0.24,0.48,0.48}{##1}}}
\@namedef{PY@tok@cpf}{\let\PY@it=\textit\def\PY@tc##1{\textcolor[rgb]{0.24,0.48,0.48}{##1}}}
\@namedef{PY@tok@c1}{\let\PY@it=\textit\def\PY@tc##1{\textcolor[rgb]{0.24,0.48,0.48}{##1}}}
\@namedef{PY@tok@cs}{\let\PY@it=\textit\def\PY@tc##1{\textcolor[rgb]{0.24,0.48,0.48}{##1}}}

\def\PYZbs{\char`\\}
\def\PYZus{\char`\_}
\def\PYZob{\char`\{}
\def\PYZcb{\char`\}}
\def\PYZca{\char`\^}
\def\PYZam{\char`\&}
\def\PYZlt{\char`\<}
\def\PYZgt{\char`\>}
\def\PYZsh{\char`\#}
\def\PYZpc{\char`\%}
\def\PYZdl{\char`\$}
\def\PYZhy{\char`\-}
\def\PYZsq{\char`\'}
\def\PYZdq{\char`\"}
\def\PYZti{\char`\~}
% for compatibility with earlier versions
\def\PYZat{@}
\def\PYZlb{[}
\def\PYZrb{]}
\makeatother


    % For linebreaks inside Verbatim environment from package fancyvrb.
    \makeatletter
        \newbox\Wrappedcontinuationbox
        \newbox\Wrappedvisiblespacebox
        \newcommand*\Wrappedvisiblespace {\textcolor{red}{\textvisiblespace}}
        \newcommand*\Wrappedcontinuationsymbol {\textcolor{red}{\llap{\tiny$\m@th\hookrightarrow$}}}
        \newcommand*\Wrappedcontinuationindent {3ex }
        \newcommand*\Wrappedafterbreak {\kern\Wrappedcontinuationindent\copy\Wrappedcontinuationbox}
        % Take advantage of the already applied Pygments mark-up to insert
        % potential linebreaks for TeX processing.
        %        {, <, #, %, $, ' and ": go to next line.
        %        _, }, ^, &, >, - and ~: stay at end of broken line.
        % Use of \textquotesingle for straight quote.
        \newcommand*\Wrappedbreaksatspecials {%
            \def\PYGZus{\discretionary{\char`\_}{\Wrappedafterbreak}{\char`\_}}%
            \def\PYGZob{\discretionary{}{\Wrappedafterbreak\char`\{}{\char`\{}}%
            \def\PYGZcb{\discretionary{\char`\}}{\Wrappedafterbreak}{\char`\}}}%
            \def\PYGZca{\discretionary{\char`\^}{\Wrappedafterbreak}{\char`\^}}%
            \def\PYGZam{\discretionary{\char`\&}{\Wrappedafterbreak}{\char`\&}}%
            \def\PYGZlt{\discretionary{}{\Wrappedafterbreak\char`\<}{\char`\<}}%
            \def\PYGZgt{\discretionary{\char`\>}{\Wrappedafterbreak}{\char`\>}}%
            \def\PYGZsh{\discretionary{}{\Wrappedafterbreak\char`\#}{\char`\#}}%
            \def\PYGZpc{\discretionary{}{\Wrappedafterbreak\char`\%}{\char`\%}}%
            \def\PYGZdl{\discretionary{}{\Wrappedafterbreak\char`\$}{\char`\$}}%
            \def\PYGZhy{\discretionary{\char`\-}{\Wrappedafterbreak}{\char`\-}}%
            \def\PYGZsq{\discretionary{}{\Wrappedafterbreak\textquotesingle}{\textquotesingle}}%
            \def\PYGZdq{\discretionary{}{\Wrappedafterbreak\char`\"}{\char`\"}}%
            \def\PYGZti{\discretionary{\char`\~}{\Wrappedafterbreak}{\char`\~}}%
        }
        % Some characters . , ; ? ! / are not pygmentized.
        % This macro makes them "active" and they will insert potential linebreaks
        \newcommand*\Wrappedbreaksatpunct {%
            \lccode`\~`\.\lowercase{\def~}{\discretionary{\hbox{\char`\.}}{\Wrappedafterbreak}{\hbox{\char`\.}}}%
            \lccode`\~`\,\lowercase{\def~}{\discretionary{\hbox{\char`\,}}{\Wrappedafterbreak}{\hbox{\char`\,}}}%
            \lccode`\~`\;\lowercase{\def~}{\discretionary{\hbox{\char`\;}}{\Wrappedafterbreak}{\hbox{\char`\;}}}%
            \lccode`\~`\:\lowercase{\def~}{\discretionary{\hbox{\char`\:}}{\Wrappedafterbreak}{\hbox{\char`\:}}}%
            \lccode`\~`\?\lowercase{\def~}{\discretionary{\hbox{\char`\?}}{\Wrappedafterbreak}{\hbox{\char`\?}}}%
            \lccode`\~`\!\lowercase{\def~}{\discretionary{\hbox{\char`\!}}{\Wrappedafterbreak}{\hbox{\char`\!}}}%
            \lccode`\~`\/\lowercase{\def~}{\discretionary{\hbox{\char`\/}}{\Wrappedafterbreak}{\hbox{\char`\/}}}%
            \catcode`\.\active
            \catcode`\,\active
            \catcode`\;\active
            \catcode`\:\active
            \catcode`\?\active
            \catcode`\!\active
            \catcode`\/\active
            \lccode`\~`\~
        }
    \makeatother

    \let\OriginalVerbatim=\Verbatim
    \makeatletter
    \renewcommand{\Verbatim}[1][1]{%
        %\parskip\z@skip
        \sbox\Wrappedcontinuationbox {\Wrappedcontinuationsymbol}%
        \sbox\Wrappedvisiblespacebox {\FV@SetupFont\Wrappedvisiblespace}%
        \def\FancyVerbFormatLine ##1{\hsize\linewidth
            \vtop{\raggedright\hyphenpenalty\z@\exhyphenpenalty\z@
                \doublehyphendemerits\z@\finalhyphendemerits\z@
                \strut ##1\strut}%
        }%
        % If the linebreak is at a space, the latter will be displayed as visible
        % space at end of first line, and a continuation symbol starts next line.
        % Stretch/shrink are however usually zero for typewriter font.
        \def\FV@Space {%
            \nobreak\hskip\z@ plus\fontdimen3\font minus\fontdimen4\font
            \discretionary{\copy\Wrappedvisiblespacebox}{\Wrappedafterbreak}
            {\kern\fontdimen2\font}%
        }%

        % Allow breaks at special characters using \PYG... macros.
        \Wrappedbreaksatspecials
        % Breaks at punctuation characters . , ; ? ! and / need catcode=\active
        \OriginalVerbatim[#1,codes*=\Wrappedbreaksatpunct]%
    }
    \makeatother

    % Exact colors from NB
    \definecolor{incolor}{HTML}{303F9F}
    \definecolor{outcolor}{HTML}{D84315}
    \definecolor{cellborder}{HTML}{CFCFCF}
    \definecolor{cellbackground}{HTML}{F7F7F7}

    % prompt
    \makeatletter
    \newcommand{\boxspacing}{\kern\kvtcb@left@rule\kern\kvtcb@boxsep}
    \makeatother
    \newcommand{\prompt}[4]{
        {\ttfamily\llap{{\color{#2}[#3]:\hspace{3pt}#4}}\vspace{-\baselineskip}}
    }
    

    
    % Prevent overflowing lines due to hard-to-break entities
    \sloppy
    % Setup hyperref package
    \hypersetup{
      breaklinks=true,  % so long urls are correctly broken across lines
      colorlinks=true,
      urlcolor=urlcolor,
      linkcolor=linkcolor,
      citecolor=citecolor,
      }
    % Slightly bigger margins than the latex defaults
    
    \geometry{verbose,tmargin=1in,bmargin=1in,lmargin=1in,rmargin=1in}
    
    

\begin{document}
    
    \maketitle
    
    

    
    \section{Postulados}\label{postulados}

\subsection{Postulado 1: Espacio de
estados}\label{postulado-1-espacio-de-estados}

Todo sistema físico aislado está asociado a un espacio de Hilbert
\(\mathcal{H}\). El estado del sistema está completamente descrito por
un \textbf{operador densidad} (matriz densidad) \(\rho\), el cual es un
operador lineal sobre \(\mathcal{H}\) que debe cumplir con tres
propiedades fundamentales. 1. Hermiticidad: \(\rho = \rho^{\dagger}\).
2. Positividad: \(\rho \ge 0\), es decir
\(\langle\psi\vert{}\rho\vert{}\psi\rangle \ge 0 \quad \forall \vert{}\psi\rangle \in \mathcal{H}\).
3. Traza unitaria: \(\operatorname{Tr}(\rho) = 1\).

Si el sistema se encuentra en un estado puro \(\vert{}\psi\rangle\),
entonces \(\rho = \vert{}\psi\rangle\langle\psi\vert{}\) (cumpliendo
\(\operatorname{Tr}(\rho^2) = 1\)). Si el sistema se encuentra en un
ensamble estadístico \(\{p_i, \vert{}\psi_i\rangle\}\) con
\(\sum_i p_i = 1\) y \(p_i \ge 0\), el operador es
\(\rho = \sum_i p_i \vert{}\psi_i\rangle\langle\psi_i\vert{}\)
(cumpliendo \(\operatorname{Tr}(\rho^2) < 1\)).

    \subsection{postulado 2: Evolucion
temporal}\label{postulado-2-evolucion-temporal}

La evolución de un sistema cuántico cerrado se describe mediante una
transformación unitaria. Si el estado del sistema en el tiempo \(t_1\)
es \(\rho(t_1)\), el estado en el tiempo \(t_2\) viene dado por:
\[\rho(t_2) = U \rho(t_1) U^\dagger\] donde \(U \equiv U(t_1, t_2)\) es
un operador unitario (\(U^\dagger U = U U^\dagger = I\)) que depende
únicamente de los tiempos \(t_1\) y \(t_2\).

\subsubsection{Ecuacion de von Neumann/Liouville-von
Neumann}\label{ecuacion-de-von-neumannliouville-von-neumann}

En tiempo continuo, si el sistema está gobernado por un Hamiltoniano
\(H(t)\), la evolución
satisface:\[\frac{d\rho(t)}{dt} = -\frac{i}{\hbar} [H(t), \rho(t)]\]donde
\([H, \rho] = H\rho - \rho H\) es el conmutador.

    \subsection{Postulado 3: Medición
cuántica}\label{postulado-3-mediciuxf3n-cuuxe1ntica}

Las mediciones cuánticas generales se describen mediante una colección
de operadores de medicion \(\{M_m\}\) que actúan sobre el espacio de
Hilbert \(\mathcal{H}\), donde el índice \(m\) denota los posibles
resultados del experimento. Estos operadores satisfacen la relación de
completítud: \[\sum_m M_m^\dagger M_m = I\] \#\#\# Probabilidad del
resultado \(m\): Data la matríz densidad \(\rho\), la probabilidad de
obtener el resultado \(m\) es
\[p(m) = \operatorname{Tr}(M_m^\dagger M_m \rho) = \operatorname{Tr}(M_m \rho M_m^\dagger)\]

\subsubsection{Estado post medición}\label{estado-post-mediciuxf3n}

Si se obtiene el resultado \(m\), el estado colapsado del sistema
inmediatamente después de la medición viene dado por: TODO: Verificar
esta formula
\[\rho' = \frac{M_m \rho M_m^\dagger}{\operatorname{Tr}(M_m^\dagger M_m \rho)} = \frac{M_m \rho M_m^\dagger}{p(m)}\]

    \subsection{postulado 4: Sistemas
compuestos}\label{postulado-4-sistemas-compuestos}

El espacio de estados de un sistema físico compuesto por múltiples
subsistemas es el producto tensorial de los espacios de Hilbert
correspondientes a cada subsistema.

Si se tienen \(n\) sistemas cuyos espacios respectivos son
\(\mathcal{H_1}, \mathcal{H_2}, ...., \mathcal{H_n}\), el espacio del
sistema global es
\[\mathcal{H} = \mathcal{H}_1 \otimes \mathcal{H}_2 \otimes \dots \otimes \mathcal{H}_n\]

Si el subsistema \(i\) se encuentra en el espacio \(\rho i\) y no
existen correlaciones ni entrelazamientos entre ellos, el estado
conjunto es:
\[\rho = \rho_1 \otimes \rho_2 \otimes \dots \otimes \rho_n\]

Para un estadp global
\(\rho_{AB} \text{ en } \mathcal{H_A} \otimes \mathcal{H_B}\), el estado
del subsistema \(A\) ( sin importar el subsistema \(B\) ) se describe
mediante la \textbf{traza parcial}
\[\rho_A = \operatorname{Tr}_B(\rho_{AB})\]

    \paragraph{Fuentes}\label{fuentes}

\begin{itemize}
\tightlist
\item
  \textbf{Nielsen M.A., \& Chuang, I.L.(2010)}. Quantum Computation and
  Quantum Information. (Capitulo 2. Introduction to quantum mechanics)
\end{itemize}

    \section{Teorema de Libertad
Unitaria}\label{teorema-de-libertad-unitaria}

Sea un operador de densidad \(\rho\). Supongamos que \(\rho\) puede ser
generado por dos ensambles diferentes de estados puros:

\begin{enumerate}
\def\labelenumi{\arabic{enumi}.}
\tightlist
\item
  Un ensamble \(E_1 = \{p_i, \vert{}\psi_i\rangle\}\), con estados (no
  necesariamente ortogonales) y probabilidades \(p_i\).
\item
  Un ensamble \(E_2 = \{q_j, \vert{}\phi_j\rangle\}\), con estados y
  probabilidades \(q_j\).
\end{enumerate}

Para simplificar la notación, definimos los ``estados ponderados'' (no
normalizados):

\(\vert{}\tilde{\psi}_i\rangle = \sqrt{p_i} \vert{}\psi_i\rangle \quad \text{y} \quad \vert{}\tilde{\phi}_j\rangle = \sqrt{q_j} \vert{}\phi_j\rangle\)

De tal forma que ambos ensambles generan la misma matriz densidad:
\[\rho = \sum_i \vert{}\tilde{\psi}_i\rangle \langle \tilde{\psi}_i\vert{} = \sum_j \vert{}\tilde{\phi}_j\rangle \langle \tilde{\phi}_j\vert{}\]

por lo tanto ambos ensambles describen el mismo operador densidad
\(\rho\) si y solo si los estados ponderados de un ensamble se pueden
expresar como combinaciones lineales de los estados del otro mediante
una matriz unitaria \(U\) (o isométrica si tienen distinto número de
elementos). Es decir:

\[\vert{}\tilde{\psi}_i\rangle = \sum_j U_{ij} \vert{}\tilde{\phi}_j\rangle\]

Si los conjuntos tienen un número distinto de vectores, se añaden
vectores nulos al conjunto más pequeño para que la matriz \(U\) sea
cuadrada y unitaria.

TODO: Validar condiciones

\subsection{\texorpdfstring{Condición Suficiente ( \(\Leftarrow\)
)}{Condición Suficiente ( \textbackslash Leftarrow )}}\label{condiciuxf3n-suficiente-leftarrow}

Supongamos que sabemos que existe una matriz unitaria \(U\) (tal que
\(U^\dagger U = I\)) que relaciona ambos ensambles:
\[\vert{}\tilde{\psi}_i\rangle = \sum_j U_{ij} \vert{}\tilde{\phi}_j\rangle\]

Debemos probar que ambos generan el mismo \(\rho\). Partimos de la
definición con los estados \(\vert{}\tilde{\psi}_i\rangle\):

\[\rho = \sum_i \vert{}\tilde{\psi}_i\rangle \langle \tilde{\psi}_i\vert{}\]

Sustituimos la relación con \(\vert{}\tilde{\phi}_j\rangle\):
\[\rho = \sum_i \left( \sum_j U_{ij} \vert{}\tilde{\phi}_j\rangle \right) \left( \sum_k U_{ik}^* \langle \tilde{\phi}_k\vert{} \right)\]

Reordenando las sumatorias (sacando la suma sobre \(i\) hacia adentro de
las constantes):
\[\rho = \sum_{j,k} \left( \sum_i U_{ik}^* U_{ij} \right) \vert{}\tilde{\phi}_j\rangle \langle \tilde{\phi}_k\vert{}\]

Dado que \(U\) es una matriz unitaria, por definición el producto de su
traspuesta conjugada \(U^\dagger\) (cuyos elementos son
\((U^\dagger)_{ki} = U_{ik}^*\)) por \(U\) es la matriz identidad \(I\).
Matemáticamente:

\[\sum_i (U^\dagger)_{ki} U_{ij} = (U^\dagger U)_{kj} = \delta_{kj}\]

\subsection{\texorpdfstring{Condición Necesaria ( \(\Rightarrow\)
)}{Condición Necesaria ( \textbackslash Rightarrow )}}\label{condiciuxf3n-necesaria-rightarrow}

Supongamos que tenemos dos ensambles que generan el mismo \(\rho\):
\[\rho = \sum_i \vert{}\tilde{\psi}_i\rangle \langle \tilde{\psi}_i\vert{} = \sum_j \vert{}\tilde{\phi}_j\rangle \langle \tilde{\phi}_j\vert{}\]

Debemos probar que existe una matriz unitaria \(U\) que los relaciona.
Como \(\rho\) es un operador hermitiano (autoadjunto) y semidefinido
positivo, podemos aplicar el teorema espectral y expresarlo en una base
ortonormal de sus autovectores \(\vert{}e_k\rangle\) con autovalores
reales y positivos \(\lambda_k > 0\):
\[\rho = \sum_k \lambda_k \vert{}e_k\rangle \langle e_k\vert{}\]

Definimos los autovectores ponderados:
\(\vert{}\tilde{e}_k\rangle = \sqrt{\lambda_k} \vert{}e_k\rangle\), por
lo que
\(\rho = \sum_k \vert{}\tilde{e}_k\rangle \langle \tilde{e}_k\vert{}\).
Cualquier estado \(\vert{}\tilde{\psi}_i\rangle\) pertenece al soporte
de \(\rho\), así que lo podemos escribir como una combinación lineal en
la base de los \(\vert{}\tilde{e}_k\rangle\):

\[\vert{}\tilde{\psi}_i\rangle = \sum_k c_{ik} \vert{}\tilde{e}_k\rangle\]

Sustituyendo esto en el operador de densidad para el primer ensamble
obtenemos:
\[\rho = \sum_i \sum_{k,l} c_{ik} c_{il}^* \vert{}\tilde{e}_k\rangle \langle \tilde{e}_l\vert{} = \sum_{k,l} \left( \sum_i c_{il}^* c_{ik} \right) \vert{}\tilde{e}_k\rangle \langle \tilde{e}_l\vert{}\]

Para que esto sea estrictamente igual a
\(\rho = \sum_k \vert{}\tilde{e}_k\rangle \langle \tilde{e}_k\vert{}\),
debe ocurrir que los coeficientes entre paréntesis se comporten como la
matriz identidad:

\[\sum_i c_{il}^* c_{ik} = \delta_{kl}\]

Esto significa que las columnas de la matriz \(C\) (formada por los
coeficientes \(c_{ik}\)) forman un conjunto de vectores ortonormales.
Toda matriz con columnas ortonormales es isométrica y se puede extender,
añadiendo filas y columnas de ceros si es necesario, a una matriz
unitaria.Aplicando exactamente el mismo razonamiento para el segundo
ensamble \(\vert{}\tilde{\phi}_j\rangle\), existe otra matriz unitaria
\(D\) (con elementos \(d_{jk}\)) tal que:

\[\vert{}\tilde{\phi}_j\rangle = \sum_k d_{jk} \vert{}\tilde{e}_k\rangle\]

Dado que ambos ensambles se pueden obtener a partir de la misma base
espectral \(\vert{}\tilde{e}_k\rangle\) mediante transformaciones
unitarias (\(C\) y \(D\)), podemos invertir algebraicamente la segunda
relación (ya que para matrices unitarias \(D^{-1} = D^\dagger\)) y
sustituirla en la primera. Esto implica que los estados
\(\vert{}\tilde{\psi}_i\rangle\) y \(\vert{}\tilde{\phi}_j\rangle\)
están directamente relacionados a través de una tercera matriz unitaria
que resulta de combinar \(C\) y \(D^\dagger\):

\[U = C D^\dagger\]

Por lo tanto, concluimos que sí existe una transformación unitaria \(U\)
que relaciona a ambos ensambles, quedando demostrada la necesidad.

    \paragraph{Fuentes}\label{fuentes}

\begin{itemize}
\tightlist
\item
  \textbf{Nielsen M.A., \& Chuang, I.L.(2010)}. Quantum Computation and
  Quantum Information. (Capitulo 2. Introduction to quantum mechanics)
\end{itemize}

    \section{Igualdad de operadores densidad para estados puros
normalizados}\label{igualdad-de-operadores-densidad-para-estados-puros-normalizados}

Consideremos dos estados cuánticos puros representados por kets
normalizados \(\vert{}\psi\rangle\) y \(\vert{}\phi\rangle\) en un
espacio de Hilbert \(\mathcal{H}\) (es decir,
\(\langle\psi\vert{}\psi\rangle = 1\) y
\(\langle\phi\vert{}\phi\rangle = 1\)). Se nos otorga como hipótesis que
sus operadores densidad asociados son idénticos:

\[\rho_\psi = \rho_\phi\]

Si dos estados puros normalizados \(\vert{}\psi\rangle\) y
\(\vert{}\phi\rangle\) poseen el mismo operador densidad
(\(\rho_\psi = \rho_\phi\)), se afirman dos propiedades fundamentales:
1. Relación Matemática: Los vectores de estado \(\vert{}\psi\rangle\) y
\(\vert{}\phi\rangle\) difieren únicamente por un factor de fase global
\(e^{i\theta}\), donde \(\theta \in \mathbb{R}\):
\[\vert{}\phi\rangle = e^{i\theta} \vert{}\psi\rangle\] 2. Significado
Físico: Ambos vectores describen exactamente el mismo estado físico. En
mecánica cuántica, la fase global no es medible ni altera ninguna
probabilidad ni valor esperado de un observable.

Para confirmarlo tenemos que.

Un estado puro se puede considerar como un ensamble de tamaño \(N = 1\)
con probabilidad \(p_1 = 1\). - Ensamble 1:
\(\{p_1 = 1, \vert{}\psi_1\rangle = \vert{}\psi\rangle\}\) - Ensamble 2:
\(\{q_1 = 1, \vert{}\phi_1\rangle = \vert{}\phi\rangle\}\)

Según el Teorema de Libertad Unitaria, ambos representan la misma matriz
de densidad \(\rho\) si y solo si existe una matriz unitaria \(U\) de
dimensiones \(1 \times 1\) que relaciona sus estados ponderados:

\[\vert{}\tilde{\psi}_1\rangle = U_{11} \vert{}\tilde{\phi}_1\rangle\]

Como \(p_1 = q_1 = 1\), los estados ponderados coinciden con los kets
originales: \(\vert{}\tilde{\psi}_1\rangle = \vert{}\psi\rangle\) y
\(\vert{}\tilde{\phi}_1\rangle = \vert{}\phi\rangle\).

Una ``matriz'' \(U\) de tamaño \(1 \times 1\) con entradas complejas que
cumple la condición de unitariedad (\(U^\dagger U = I_{1\times1}\)) no
es más que un número complejo de módulo \(1\), de la forma
\(u_{11} = e^{-i\theta}\). Sustituyendo

\[\vert{}\psi\rangle = e^{-i\theta} \vert{}\phi\rangle \implies \vert{}\phi\rangle = e^{i\theta} \vert{}\psi\rangle\]

\paragraph{Fuentes}\label{fuentes}

\begin{itemize}
\tightlist
\item
  \textbf{Nielsen M.A., \& Chuang, I.L.(2010)}. Quantum Computation and
  Quantum Information. (Capitulo 2.4. The density operator)
\end{itemize}

    \section{Cálculo y comparación de Matrices de
Densidad.}\label{cuxe1lculo-y-comparaciuxf3n-de-matrices-de-densidad.}

El operador densidad \(\rho\) para un ensamble de estados puros
\(\{p_k, \vert{}\psi_k\rangle\}\) se define como la suma ponderada de
los proyectores de cada estado:
\[\rho = \sum_k p_k \vert{}\psi_k\rangle\langle\psi_k\vert{}\] Para
encontrar la representación matricial, calcularemos el producto exterior
\(\vert{}\psi_k\rangle\langle\psi_k\vert{}\) para cada estado, lo
multiplicaremos por su probabilidad asociada \(p_k\), y sumaremos los
resultados.

    \subsection{Matriz de densidad del Conjunto
A}\label{matriz-de-densidad-del-conjunto-a}

A1.
\(\vert{}\psi_{1}\rangle = \frac{3}{\sqrt{10}}\vert{}0\rangle + \frac{1}{\sqrt{10}}\vert{}1\rangle\),
con \(p_{1} = \frac{2}{5}\)

A2.
\(\vert{}\psi_{2}\rangle = \frac{2}{\sqrt{13}}\vert{}0\rangle - \frac{3i}{\sqrt{13}}\vert{}1\rangle\),
con \(p_{2} = \frac{3}{10}\)

A3.
\(\vert{}\psi_{3}\rangle = \frac{i}{2}\vert{}0\rangle + \frac{\sqrt{3}}{2}\vert{}1\rangle\),
con \(p_{3} = \frac{3}{10}\)

    \begin{tcolorbox}[breakable, size=fbox, boxrule=1pt, pad at break*=1mm,colback=cellbackground, colframe=cellborder]
\prompt{In}{incolor}{32}{\boxspacing}
\begin{Verbatim}[commandchars=\\\{\}]
\PY{c+c1}{\PYZsh{} Implementacion usando Qiskit}
\PY{k+kn}{import}\PY{+w}{ }\PY{n+nn}{numpy}\PY{+w}{ }\PY{k}{as}\PY{+w}{ }\PY{n+nn}{np}
\PY{k+kn}{from}\PY{+w}{ }\PY{n+nn}{qiskit}\PY{n+nn}{.}\PY{n+nn}{quantum\PYZus{}info}\PY{+w}{ }\PY{k+kn}{import} \PY{n}{Statevector}\PY{p}{,} \PY{n}{DensityMatrix}\PY{p}{,} \PY{n}{partial\PYZus{}trace}
\PY{k+kn}{from}\PY{+w}{ }\PY{n+nn}{IPython}\PY{n+nn}{.}\PY{n+nn}{display}\PY{+w}{ }\PY{k+kn}{import} \PY{n}{display}\PY{p}{,} \PY{n}{Math}

\PY{c+c1}{\PYZsh{} Importamos librerias de tipado de datos}
\PY{k+kn}{import}\PY{+w}{ }\PY{n+nn}{numpy}\PY{n+nn}{.}\PY{n+nn}{typing}\PY{+w}{ }\PY{k}{as}\PY{+w}{ }\PY{n+nn}{npt}

\PY{k}{def}\PY{+w}{ }\PY{n+nf}{generar\PYZus{}conjunto}\PY{p}{(}
    \PY{n}{conjunto}\PY{p}{:} \PY{n}{List}\PY{p}{[}\PY{n}{Dict}\PY{p}{[}\PY{l+s+s2}{\PYZdq{}}\PY{l+s+s2}{p}\PY{l+s+s2}{\PYZdq{}}\PY{o}{|}\PY{l+s+s2}{\PYZdq{}}\PY{l+s+s2}{kets}\PY{l+s+s2}{\PYZdq{}}\PY{p}{,} \PY{n}{Any}\PY{p}{]}\PY{p}{]}\PY{p}{,}
    \PY{n}{nombre}
\PY{p}{)} \PY{o}{\PYZhy{}}\PY{o}{\PYZgt{}} \PY{n}{DensityMatrix}\PY{p}{:}
\PY{+w}{    }\PY{l+s+sd}{\PYZdq{}\PYZdq{}\PYZdq{} Esta funcion genera todo el calculo del conjunto y lo regresa en un solo valor\PYZdq{}\PYZdq{}\PYZdq{}}
    \PY{n}{md\PYZus{}conjunto} \PY{o}{=} \PY{k+kc}{None}
    \PY{k}{for} \PY{n}{item} \PY{o+ow}{in} \PY{n}{conjunto}\PY{p}{:}
        \PY{n}{p} \PY{o}{=} \PY{n}{item}\PY{p}{[}\PY{l+s+s1}{\PYZsq{}}\PY{l+s+s1}{p}\PY{l+s+s1}{\PYZsq{}}\PY{p}{]}
        \PY{n}{kets} \PY{o}{=} \PY{n}{item}\PY{p}{[}\PY{l+s+s1}{\PYZsq{}}\PY{l+s+s1}{kets}\PY{l+s+s1}{\PYZsq{}}\PY{p}{]}
        
        \PY{n}{estado} \PY{o}{=} \PY{n}{Statevector}\PY{p}{(}\PY{n}{kets}\PY{p}{)}
        \PY{n}{md} \PY{o}{=} \PY{n}{DensityMatrix}\PY{p}{(}\PY{n}{estado}\PY{p}{)}

        \PY{k}{if} \PY{n}{md\PYZus{}conjunto} \PY{o+ow}{is} \PY{k+kc}{None}\PY{p}{:}
            \PY{n}{md\PYZus{}conjunto} \PY{o}{=} \PY{p}{(}\PY{n}{p} \PY{o}{*} \PY{n}{md}\PY{p}{)}
        \PY{k}{else}\PY{p}{:}
            \PY{n}{md\PYZus{}conjunto} \PY{o}{+}\PY{o}{=} \PY{p}{(}\PY{n}{p} \PY{o}{*} \PY{n}{md}\PY{p}{)}

    \PY{n}{traza} \PY{o}{=} \PY{n}{np}\PY{o}{.}\PY{n}{trace}\PY{p}{(}\PY{n}{md\PYZus{}conjunto}\PY{p}{)}
    \PY{n}{es\PYZus{}hermetica} \PY{o}{=} \PY{n}{np}\PY{o}{.}\PY{n}{allclose}\PY{p}{(}\PY{n}{md\PYZus{}conjunto}\PY{o}{.}\PY{n}{data}\PY{p}{,} \PY{n}{md\PYZus{}conjunto}\PY{o}{.}\PY{n}{data}\PY{o}{.}\PY{n}{conj}\PY{p}{(}\PY{p}{)}\PY{o}{.}\PY{n}{T}\PY{p}{)}
    
    \PY{n}{display}\PY{p}{(}\PY{n}{Math}\PY{p}{(}\PY{l+s+sa}{rf}\PY{l+s+s2}{\PYZdq{}}\PY{l+s+s2}{\PYZbs{}}\PY{l+s+s2}{text}\PY{l+s+se}{\PYZob{}\PYZob{}}\PY{l+s+s2}{ Matriz de Densidad del Conjunto: }\PY{l+s+se}{\PYZcb{}\PYZcb{}}\PY{l+s+s2}{ }\PY{l+s+si}{\PYZob{}}\PY{n}{nombre}\PY{l+s+si}{\PYZcb{}}\PY{l+s+s2}{\PYZdq{}}\PY{p}{)}\PY{p}{)}
    \PY{n}{display}\PY{p}{(}\PY{n}{md\PYZus{}conjunto}\PY{o}{.}\PY{n}{draw}\PY{p}{(}\PY{l+s+s1}{\PYZsq{}}\PY{l+s+s1}{latex}\PY{l+s+s1}{\PYZsq{}}\PY{p}{,}  \PY{n}{prefix}\PY{o}{=}\PY{l+s+sa}{f}\PY{l+s+s2}{\PYZdq{}}\PY{l+s+si}{\PYZob{}}\PY{n}{nombre}\PY{l+s+si}{\PYZcb{}}\PY{l+s+s2}{ = }\PY{l+s+s2}{\PYZdq{}}\PY{p}{)}\PY{p}{)}
    \PY{n}{display}\PY{p}{(}\PY{n}{Math}\PY{p}{(}\PY{l+s+sa}{rf}\PY{l+s+s2}{\PYZdq{}}\PY{l+s+s2}{\PYZbs{}}\PY{l+s+s2}{text}\PY{l+s+se}{\PYZob{}\PYZob{}}\PY{l+s+s2}{Tr}\PY{l+s+se}{\PYZcb{}\PYZcb{}}\PY{l+s+s2}{(}\PY{l+s+s2}{\PYZbs{}}\PY{l+s+s2}{rho\PYZus{}}\PY{l+s+si}{\PYZob{}}\PY{n}{nombre}\PY{l+s+si}{\PYZcb{}}\PY{l+s+s2}{) = }\PY{l+s+si}{\PYZob{}}\PY{n}{np}\PY{o}{.}\PY{n}{round}\PY{p}{(}\PY{n}{traza}\PY{o}{.}\PY{n}{real}\PY{p}{,}\PY{+w}{ }\PY{l+m+mi}{4}\PY{p}{)}\PY{l+s+si}{\PYZcb{}}\PY{l+s+s2}{\PYZdq{}}\PY{p}{)}\PY{p}{)}
    \PY{n}{display}\PY{p}{(}\PY{n}{Math}\PY{p}{(}\PY{l+s+sa}{rf}\PY{l+s+s2}{\PYZdq{}}\PY{l+s+s2}{\PYZbs{}}\PY{l+s+s2}{text}\PY{l+s+se}{\PYZob{}\PYZob{}}\PY{l+s+s2}{¿Es Hermítica?}\PY{l+s+se}{\PYZcb{}\PYZcb{}}\PY{l+s+s2}{: }\PY{l+s+si}{\PYZob{}}\PY{n}{es\PYZus{}hermetica}\PY{l+s+si}{\PYZcb{}}\PY{l+s+s2}{\PYZdq{}}\PY{p}{)}\PY{p}{)}
    \PY{k}{return} \PY{n}{md\PYZus{}conjunto}
\end{Verbatim}
\end{tcolorbox}

    \begin{tcolorbox}[breakable, size=fbox, boxrule=1pt, pad at break*=1mm,colback=cellbackground, colframe=cellborder]
\prompt{In}{incolor}{33}{\boxspacing}
\begin{Verbatim}[commandchars=\\\{\}]
\PY{n}{conjunto\PYZus{}A} \PY{o}{=} \PY{n}{generar\PYZus{}conjunto}\PY{p}{(}
    \PY{p}{[}
        \PY{p}{\PYZob{}}\PY{l+s+s2}{\PYZdq{}}\PY{l+s+s2}{p}\PY{l+s+s2}{\PYZdq{}}\PY{p}{:} \PY{l+m+mi}{2}\PY{o}{/}\PY{l+m+mi}{5}\PY{p}{,} \PY{l+s+s2}{\PYZdq{}}\PY{l+s+s2}{kets}\PY{l+s+s2}{\PYZdq{}}\PY{p}{:} \PY{p}{[}\PY{l+m+mi}{3}\PY{o}{/}\PY{n}{np}\PY{o}{.}\PY{n}{sqrt}\PY{p}{(}\PY{l+m+mi}{10}\PY{p}{)}\PY{p}{,} \PY{l+m+mi}{1}\PY{o}{/}\PY{n}{np}\PY{o}{.}\PY{n}{sqrt}\PY{p}{(}\PY{l+m+mi}{10}\PY{p}{)}\PY{p}{]}\PY{p}{\PYZcb{}}\PY{p}{,}
        \PY{p}{\PYZob{}}\PY{l+s+s2}{\PYZdq{}}\PY{l+s+s2}{p}\PY{l+s+s2}{\PYZdq{}}\PY{p}{:} \PY{l+m+mi}{3}\PY{o}{/}\PY{l+m+mi}{10}\PY{p}{,} \PY{l+s+s2}{\PYZdq{}}\PY{l+s+s2}{kets}\PY{l+s+s2}{\PYZdq{}}\PY{p}{:} \PY{p}{[}\PY{l+m+mi}{2}\PY{o}{/}\PY{n}{np}\PY{o}{.}\PY{n}{sqrt}\PY{p}{(}\PY{l+m+mi}{13}\PY{p}{)}\PY{p}{,} \PY{o}{\PYZhy{}}\PY{l+m+mi}{3}\PY{n}{j}\PY{o}{/}\PY{n}{np}\PY{o}{.}\PY{n}{sqrt}\PY{p}{(}\PY{l+m+mi}{13}\PY{p}{)}\PY{p}{]}\PY{p}{\PYZcb{}}\PY{p}{,}
        \PY{p}{\PYZob{}}\PY{l+s+s2}{\PYZdq{}}\PY{l+s+s2}{p}\PY{l+s+s2}{\PYZdq{}}\PY{p}{:} \PY{l+m+mi}{3}\PY{o}{/}\PY{l+m+mi}{10}\PY{p}{,} \PY{l+s+s2}{\PYZdq{}}\PY{l+s+s2}{kets}\PY{l+s+s2}{\PYZdq{}}\PY{p}{:} \PY{p}{[}\PY{l+m+mi}{1}\PY{n}{j}\PY{o}{/}\PY{l+m+mi}{2}\PY{p}{,} \PY{n}{np}\PY{o}{.}\PY{n}{sqrt}\PY{p}{(}\PY{l+m+mi}{3}\PY{p}{)}\PY{o}{/}\PY{l+m+mi}{2}\PY{p}{]}\PY{p}{\PYZcb{}}
    \PY{p}{]}\PY{p}{,}
    \PY{l+s+s2}{\PYZdq{}}\PY{l+s+s2}{A}\PY{l+s+s2}{\PYZdq{}}
\PY{p}{)}
\end{Verbatim}
\end{tcolorbox}

    $\displaystyle \text{ Matriz de Densidad del Conjunto: } A$

    
    $$
A = 
\begin{bmatrix}
0.5273076923 & 0.12 + 0.268365349 i  \\
 0.12 - 0.268365349 i & 0.4726923077  \\
 \end{bmatrix}
$$

    
    $\displaystyle \text{Tr}(\rho_A) = 1.0$

    
    $\displaystyle \text{¿Es Hermítica?}: True$

    
    \subsection{Matriz de densidad del Conjunto
B}\label{matriz-de-densidad-del-conjunto-b}

B1. \(\vert{}\psi_{1}\rangle = \vert{0}\rangle\), con
\(p_1 = \frac{1}{2}\)

B2.
\(\vert\psi_{2}\rangle = \frac{1}{\sqrt{2}}\vert{0}\rangle + \frac{1}{\sqrt{2}}\vert{1}\rangle\),
con \(p_2 = \frac{1}{3}\)

B3.
\(\vert\psi_{3}\rangle = \frac{1}{\sqrt{2}}\vert{0}\rangle - \frac{i}{\sqrt{2}}\vert{1}\rangle\),
con \(p_3 = \frac{1}{6}\)

    \begin{tcolorbox}[breakable, size=fbox, boxrule=1pt, pad at break*=1mm,colback=cellbackground, colframe=cellborder]
\prompt{In}{incolor}{35}{\boxspacing}
\begin{Verbatim}[commandchars=\\\{\}]
\PY{n}{conjunto\PYZus{}B} \PY{o}{=} \PY{n}{generar\PYZus{}conjunto}\PY{p}{(}
    \PY{p}{[}
        \PY{p}{\PYZob{}}\PY{l+s+s2}{\PYZdq{}}\PY{l+s+s2}{p}\PY{l+s+s2}{\PYZdq{}}\PY{p}{:} \PY{l+m+mi}{1}\PY{o}{/}\PY{l+m+mi}{2}\PY{p}{,} \PY{l+s+s2}{\PYZdq{}}\PY{l+s+s2}{kets}\PY{l+s+s2}{\PYZdq{}}\PY{p}{:} \PY{p}{[}\PY{l+m+mi}{1}\PY{p}{,}\PY{l+m+mi}{0}\PY{p}{]}\PY{p}{\PYZcb{}}\PY{p}{,}
        \PY{p}{\PYZob{}}\PY{l+s+s2}{\PYZdq{}}\PY{l+s+s2}{p}\PY{l+s+s2}{\PYZdq{}}\PY{p}{:} \PY{l+m+mi}{1}\PY{o}{/}\PY{l+m+mi}{3}\PY{p}{,} \PY{l+s+s2}{\PYZdq{}}\PY{l+s+s2}{kets}\PY{l+s+s2}{\PYZdq{}}\PY{p}{:} \PY{p}{[}\PY{l+m+mi}{1}\PY{o}{/}\PY{n}{np}\PY{o}{.}\PY{n}{sqrt}\PY{p}{(}\PY{l+m+mi}{2}\PY{p}{)}\PY{p}{,} \PY{l+m+mi}{1}\PY{o}{/}\PY{n}{np}\PY{o}{.}\PY{n}{sqrt}\PY{p}{(}\PY{l+m+mi}{2}\PY{p}{)}\PY{p}{]}\PY{p}{\PYZcb{}}\PY{p}{,}
        \PY{p}{\PYZob{}}\PY{l+s+s2}{\PYZdq{}}\PY{l+s+s2}{p}\PY{l+s+s2}{\PYZdq{}}\PY{p}{:} \PY{l+m+mi}{1}\PY{o}{/}\PY{l+m+mi}{6}\PY{p}{,} \PY{l+s+s2}{\PYZdq{}}\PY{l+s+s2}{kets}\PY{l+s+s2}{\PYZdq{}}\PY{p}{:} \PY{p}{[}\PY{l+m+mi}{1}\PY{o}{/}\PY{n}{np}\PY{o}{.}\PY{n}{sqrt}\PY{p}{(}\PY{l+m+mi}{2}\PY{p}{)}\PY{p}{,} \PY{o}{\PYZhy{}}\PY{l+m+mi}{1}\PY{n}{j}\PY{o}{/}\PY{n}{np}\PY{o}{.}\PY{n}{sqrt}\PY{p}{(}\PY{l+m+mi}{2}\PY{p}{)}\PY{p}{]}\PY{p}{\PYZcb{}}
    \PY{p}{]}\PY{p}{,}
    \PY{l+s+s2}{\PYZdq{}}\PY{l+s+s2}{B}\PY{l+s+s2}{\PYZdq{}}
\PY{p}{)}
\end{Verbatim}
\end{tcolorbox}

    $\displaystyle \text{ Matriz de Densidad del Conjunto: } B$

    
    $$
B = 
\begin{bmatrix}
\frac{3}{4} & \frac{1}{6} + \frac{i}{12}  \\
 \frac{1}{6} - \frac{i}{12} & \frac{1}{4}  \\
 \end{bmatrix}
$$

    
    $\displaystyle \text{Tr}(\rho_B) = 1.0$

    
    $\displaystyle \text{¿Es Hermítica?}: True$

    
    \begin{tcolorbox}[breakable, size=fbox, boxrule=1pt, pad at break*=1mm,colback=cellbackground, colframe=cellborder]
\prompt{In}{incolor}{36}{\boxspacing}
\begin{Verbatim}[commandchars=\\\{\}]
\PY{c+c1}{\PYZsh{}\PYZsh{} Comprobacion de matrices}
\PY{n+nb}{print}\PY{p}{(}\PY{l+s+s2}{\PYZdq{}}\PY{l+s+s2}{Las Matrices densidad de los conjuntos A y B son iguales?}\PY{l+s+s2}{\PYZdq{}}\PY{p}{,} \PY{n}{np}\PY{o}{.}\PY{n}{allclose}\PY{p}{(}\PY{n}{conjunto\PYZus{}A}\PY{o}{.}\PY{n}{data}\PY{p}{,} \PY{n}{conjunto\PYZus{}B}\PY{o}{.}\PY{n}{data}\PY{p}{)}\PY{p}{)}
\end{Verbatim}
\end{tcolorbox}

    \begin{Verbatim}[commandchars=\\\{\}]
Las Matrices densidad de los conjuntos A y B son iguales? False
    \end{Verbatim}

    Al comparar las matrices elemento por elemento vemos que los valores no
coinciden, tomemos de ejemplo la probabilidadde medir el estado
\(\vert{0}\rangle\) en la base computacional, para el \textbf{conjunto
A} tenemos que la probabilidad es de \textbf{25.73\%} mientras que para
el \textbf{conjunto B tenemos} el \textbf{75\%} \(\frac{3}{4}\), por lo
que tenemos que \(\rho_A \neq \rho_B\), el Teorema de la Libertad
Unitaria \textbf{no existe} ninguna tranformación unitaria que relacione
los estados de ambos ensambles.

    \section{El Estado de Bell y su Matriz de Densidad
Global}\label{el-estado-de-bell-y-su-matriz-de-densidad-global}

Consideremos un sistema bipartito compuesto por dos qubits (A y B) en el
estado entrelazado:
\[\vert{}\Phi^+\rangle = \frac{1}{\sqrt{2}}(\vert{}00\rangle + \vert{}11\rangle)\]
El operador densidad para este estado puro del sistema global es el
proyector: \[\rho_{AB} = \vert{}\Phi^+\rangle \langle \Phi^+\vert{}\]
Sustituyendo el ket y el bra correspondientes:
\[\rho_{AB} = \left( \frac{1}{\sqrt{2}}\vert{}00\rangle + \frac{1}{\sqrt{2}}\vert{}11\rangle \right) \left( \frac{1}{\sqrt{2}}\langle 00\vert{} + \frac{1}{\sqrt{2}}\langle 11\vert{} \right)\]
Haciendo la multiplicación distributiva (producto exterior):
\[\rho_{AB} = \frac{1}{2} \left( \vert{}00\rangle\langle 00\vert{} + \vert{}00\rangle\langle 11\vert{} + \vert{}11\rangle\langle 00\vert{} + \vert{}11\rangle\langle 11\vert{} \right)\]

En su representación matricial (en la base computacional
\(\vert{}00\rangle, \vert{}01\rangle, \vert{}10\rangle, \vert{}11\rangle\)),
obtenemos una matriz de \(4 \times 4\):

\[\rho_{AB} = \frac{1}{2} \begin{pmatrix} 1 & 0 & 0 & 1 \\ 0 & 0 & 0 & 0 \\ 0 & 0 & 0 & 0 \\ 1 & 0 & 0 & 1 \end{pmatrix} = \begin{pmatrix} 0.5 & 0 & 0 & 0.5 \\ 0 & 0 & 0 & 0 \\ 0 & 0 & 0 & 0 \\ 0.5 & 0 & 0 & 0.5 \end{pmatrix}\]

\(\text{Tr}(\rho_{AB}) = 0.5 + 0.5 = 1\) y es una matriz Hermítica, por
lo que cumple perfectamente los postulados para un sistema válido

\subsection{\texorpdfstring{La Traza Parcial (Matriz de Densidad
Reducida
\(\rho_A\))}{La Traza Parcial (Matriz de Densidad Reducida \textbackslash rho\_A)}}\label{la-traza-parcial-matriz-de-densidad-reducida-rho_a}

Recordemos que los estados de la base son en realidad productos
tensoriales
(\(\vert{}00\rangle = \vert{}0_A\rangle \otimes \vert{}0_B\rangle\)). Al
aplicar los bras y kets del Qubit B, solo afectan a la parte B del
producto:Evaluando con \(\langle 0_B\vert{}\) y
\(\vert{} 0_B \rangle\):Solo sobrevive el término que tiene
\(\vert{}0_B\rangle\langle 0_B\vert{}\) en el centro, es decir, la parte
de \(\vert{}00\rangle\langle 00\vert{}\) que deja a
\(\vert{}0_A\rangle\langle 0_A\vert{}\). Los términos cruzados como
\(\vert{}00\rangle\langle 11\vert{}\) se anulan porque
\(\langle 0_B\vert{}1_B\rangle = 0\) (son ortogonales).

\[\langle 0_B \vert{} \rho_{AB} \vert{} 0_B \rangle = \frac{1}{2} \vert{}0_A\rangle\langle 0_A\vert{}\]

Evaluando con \(\langle 1_B\vert{}\) y \(\vert{} 1_B \rangle\): De forma
análoga, solo sobrevive el término de
\(\vert{}11\rangle\langle 11\vert{}\), dejando:

\[\rho_A = \frac{1}{2} \vert{}0_A\rangle\langle 0_A\vert{} + \frac{1}{2} \vert{}1_A\rangle\langle 1_A\vert{}\]

En representación matricial (\(2 \times 2\)):

\[\rho_A = \begin{pmatrix} \frac{1}{2} & 0 \\ 0 & \frac{1}{2} \end{pmatrix} = \begin{pmatrix} 0.5 & 0 \\ 0 & 0.5 \end{pmatrix} = \frac{1}{2} I\]

    \begin{tcolorbox}[breakable, size=fbox, boxrule=1pt, pad at break*=1mm,colback=cellbackground, colframe=cellborder]
\prompt{In}{incolor}{37}{\boxspacing}
\begin{Verbatim}[commandchars=\\\{\}]
\PY{n}{eb} \PY{o}{=} \PY{n}{Statevector}\PY{p}{(}\PY{p}{[}\PY{l+m+mi}{1}\PY{o}{/}\PY{n}{np}\PY{o}{.}\PY{n}{sqrt}\PY{p}{(}\PY{l+m+mi}{2}\PY{p}{)}\PY{p}{,} \PY{l+m+mi}{0}\PY{p}{,} \PY{l+m+mi}{0}\PY{p}{,} \PY{l+m+mi}{1}\PY{o}{/}\PY{n}{np}\PY{o}{.}\PY{n}{sqrt}\PY{p}{(}\PY{l+m+mi}{2}\PY{p}{)}\PY{p}{]}\PY{p}{)}

\PY{n}{md\PYZus{}ab} \PY{o}{=} \PY{n}{DensityMatrix}\PY{p}{(}\PY{n}{eb}\PY{p}{)}
\PY{n}{display}\PY{p}{(}\PY{n}{Math}\PY{p}{(}\PY{l+s+sa}{rf}\PY{l+s+s2}{\PYZdq{}}\PY{l+s+s2}{\PYZbs{}}\PY{l+s+s2}{text}\PY{l+s+se}{\PYZob{}\PYZob{}}\PY{l+s+s2}{Matriz de densidad global del estado de Bell }\PY{l+s+se}{\PYZcb{}\PYZcb{}}\PY{l+s+s2}{ md\PYZus{}}\PY{l+s+se}{\PYZob{}\PYZob{}}\PY{l+s+s2}{ab}\PY{l+s+se}{\PYZcb{}\PYZcb{}}\PY{l+s+s2}{\PYZdq{}}\PY{p}{)}\PY{p}{)}
\PY{n}{display}\PY{p}{(}\PY{n}{md\PYZus{}ab}\PY{o}{.}\PY{n}{draw}\PY{p}{(}\PY{l+s+s1}{\PYZsq{}}\PY{l+s+s1}{latex}\PY{l+s+s1}{\PYZsq{}}\PY{p}{,} \PY{n}{prefix}\PY{o}{=}\PY{l+s+s2}{\PYZdq{}}\PY{l+s+se}{\PYZbs{}\PYZbs{}}\PY{l+s+s2}{rho = }\PY{l+s+s2}{\PYZdq{}}\PY{p}{)}\PY{p}{)}
\PY{n}{md\PYZus{}a} \PY{o}{=} \PY{n}{partial\PYZus{}trace}\PY{p}{(}\PY{n}{md\PYZus{}ab}\PY{p}{,} \PY{p}{[}\PY{l+m+mi}{0}\PY{p}{]}\PY{p}{)}
\PY{n}{display}\PY{p}{(}\PY{n}{Math}\PY{p}{(}\PY{l+s+sa}{rf}\PY{l+s+s2}{\PYZdq{}}\PY{l+s+s2}{\PYZbs{}}\PY{l+s+s2}{text}\PY{l+s+se}{\PYZob{}\PYZob{}}\PY{l+s+s2}{Matriz de densidad reducida }\PY{l+s+se}{\PYZcb{}\PYZcb{}}\PY{l+s+s2}{ }\PY{l+s+s2}{\PYZbs{}}\PY{l+s+s2}{rho\PYZus{}A }\PY{l+s+s2}{\PYZbs{}}\PY{l+s+s2}{text}\PY{l+s+se}{\PYZob{}\PYZob{}}\PY{l+s+s2}{ \PYZhy{} Estado Máximamente Mixto:}\PY{l+s+se}{\PYZcb{}\PYZcb{}}\PY{l+s+s2}{\PYZdq{}}\PY{p}{)}\PY{p}{)}
\PY{n}{display}\PY{p}{(}\PY{n}{md\PYZus{}a}\PY{o}{.}\PY{n}{draw}\PY{p}{(}\PY{l+s+s1}{\PYZsq{}}\PY{l+s+s1}{latex}\PY{l+s+s1}{\PYZsq{}}\PY{p}{,} \PY{n}{prefix}\PY{o}{=}\PY{l+s+s2}{\PYZdq{}}\PY{l+s+se}{\PYZbs{}\PYZbs{}}\PY{l+s+s2}{rho\PYZus{}A = }\PY{l+s+s2}{\PYZdq{}}\PY{p}{)}\PY{p}{)}
\end{Verbatim}
\end{tcolorbox}

    $\displaystyle \text{Matriz de densidad global del estado de Bell } md_{ab}$

    
    $$
\rho = 
\begin{bmatrix}
\frac{1}{2} & 0 & 0 & \frac{1}{2}  \\
 0 & 0 & 0 & 0  \\
 0 & 0 & 0 & 0  \\
 \frac{1}{2} & 0 & 0 & \frac{1}{2}  \\
 \end{bmatrix}
$$

    
    $\displaystyle \text{Matriz de densidad reducida } \rho_A \text{ - Estado Máximamente Mixto:}$

    
    $$
\rho_A = 
\begin{bmatrix}
\frac{1}{2} & 0  \\
 0 & \frac{1}{2}  \\
 \end{bmatrix}
$$

    
    \paragraph{Fuentes}\label{fuentes}

\begin{itemize}
\tightlist
\item
  \textbf{Nielsen M.A., \& Chuang, I.L.(2010)}. Quantum Computation and
  Quantum Information. (Capitulo 2.4.3 The reduced density operator)
\end{itemize}

    \section{Demostración de la Hermiticidad del Operador
Densidad}\label{demostraciuxf3n-de-la-hermiticidad-del-operador-densidad}

Para demostrar que \(\rho\) es Hermítico, debemos calcular
\(\rho^\dagger\) y comprobar que el resultado es igual a
\(\rho\).Tomamos el conjugado traspuesto de toda la expresión del
operador densidad:

\[\rho^\dagger = \left( \sum_i p_i \vert{}\psi_i\rangle \langle \psi_i\vert{} \right)^\dagger\]

Tomando en cuenta la matriz densidad del estado de Bell y la matriz
densidad reducida

    \begin{tcolorbox}[breakable, size=fbox, boxrule=1pt, pad at break*=1mm,colback=cellbackground, colframe=cellborder]
\prompt{In}{incolor}{38}{\boxspacing}
\begin{Verbatim}[commandchars=\\\{\}]
\PY{k}{def}\PY{+w}{ }\PY{n+nf}{comprobar\PYZus{}hermiticidad}\PY{p}{(}\PY{n}{md}\PY{p}{,} \PY{n}{nombre}\PY{p}{)}\PY{p}{:}
    \PY{n}{matriz} \PY{o}{=} \PY{n}{md}\PY{o}{.}\PY{n}{data}
    \PY{n}{matriz\PYZus{}daga} \PY{o}{=} \PY{n}{matriz}\PY{o}{.}\PY{n}{conj}\PY{p}{(}\PY{p}{)}\PY{o}{.}\PY{n}{T}

    \PY{n}{es\PYZus{}hermitica} \PY{o}{=} \PY{n}{np}\PY{o}{.}\PY{n}{allclose}\PY{p}{(}\PY{n}{matriz}\PY{p}{,} \PY{n}{matriz\PYZus{}daga}\PY{p}{)}
    \PY{n}{display}\PY{p}{(}\PY{n}{Math}\PY{p}{(}\PY{l+s+sa}{rf}\PY{l+s+s2}{\PYZdq{}}\PY{l+s+s2}{\PYZbs{}}\PY{l+s+s2}{text}\PY{l+s+se}{\PYZob{}\PYZob{}}\PY{l+s+s2}{Verificando: }\PY{l+s+se}{\PYZcb{}\PYZcb{}}\PY{l+s+s2}{ }\PY{l+s+si}{\PYZob{}}\PY{n}{nombre}\PY{l+s+si}{\PYZcb{}}\PY{l+s+s2}{\PYZdq{}}\PY{p}{)}\PY{p}{)}
    \PY{n}{display}\PY{p}{(}\PY{n}{Math}\PY{p}{(}\PY{l+s+sa}{rf}\PY{l+s+s2}{\PYZdq{}}\PY{l+s+s2}{\PYZbs{}}\PY{l+s+s2}{text}\PY{l+s+se}{\PYZob{}\PYZob{}}\PY{l+s+s2}{¿Es estrictamente igual a su traspuesta conjugada? \PYZhy{}\PYZgt{} }\PY{l+s+si}{\PYZob{}}\PY{n}{es\PYZus{}hermitica}\PY{l+s+si}{\PYZcb{}}\PY{l+s+se}{\PYZcb{}\PYZcb{}}\PY{l+s+s2}{\PYZdq{}}\PY{p}{)}\PY{p}{)}

\PY{n}{comprobar\PYZus{}hermiticidad}\PY{p}{(}\PY{n}{md\PYZus{}ab}\PY{p}{,} \PY{l+s+sa}{f}\PY{l+s+s2}{\PYZdq{}}\PY{l+s+se}{\PYZbs{}\PYZbs{}}\PY{l+s+s2}{rho\PYZus{}}\PY{l+s+se}{\PYZob{}\PYZob{}}\PY{l+s+s2}{AB}\PY{l+s+se}{\PYZcb{}\PYZcb{}}\PY{l+s+s2}{\PYZdq{}}\PY{p}{)}
\PY{n}{comprobar\PYZus{}hermiticidad}\PY{p}{(}\PY{n}{md\PYZus{}a}\PY{p}{,} \PY{l+s+sa}{f}\PY{l+s+s2}{\PYZdq{}}\PY{l+s+se}{\PYZbs{}\PYZbs{}}\PY{l+s+s2}{rho\PYZus{}}\PY{l+s+se}{\PYZob{}\PYZob{}}\PY{l+s+s2}{A}\PY{l+s+se}{\PYZcb{}\PYZcb{}}\PY{l+s+s2}{\PYZdq{}}\PY{p}{)}
\end{Verbatim}
\end{tcolorbox}

    $\displaystyle \text{Verificando: } \rho_{AB}$

    
    $\displaystyle \text{¿Es estrictamente igual a su traspuesta conjugada? -> True}$

    
    $\displaystyle \text{Verificando: } \rho_{A}$

    
    $\displaystyle \text{¿Es estrictamente igual a su traspuesta conjugada? -> True}$

    
    Se aplicó la operación de transpuesta conjugada. En ambos escenarios, la
función \texttt{np.allclose()} confirmó una igualdad absoluta entre
\(\rho\) y \(\rho^\dagger\), comprobando empíricamente que la
hermiticidad se preserva sin importar si el sistema se encuentra en un
estado puro o en una mezcla estadística, validando así la robustez del
modelo teórico


    % Add a bibliography block to the postdoc
    
    
    
\end{document}
